% Generated draft: requires mathematical and source review.
\documentclass[11pt]{article}
\usepackage[margin=25mm]{geometry}
\usepackage{fontspec}
\setmainfont{Latin Modern Roman}
\usepackage{amsmath,amssymb,mathtools,unicode-math}
\setmathfont{Latin Modern Math}
\usepackage{enumitem,needspace,xurl,hyperref}
\setlength{\parindent}{0pt}
\setlength{\parskip}{5pt}
\setlength{\emergencystretch}{3em}
\newcommand{\N}{\mathbb{N}}
\newcommand{\Z}{\mathbb{Z}}
\newcommand{\Q}{\mathbb{Q}}
\newcommand{\R}{\mathbb{R}}
\newcommand{\C}{\mathbb{C}}
\newcommand{\F}{\mathbb{F}}
\newcommand{\A}{\mathbb{A}}
\newcommand{\PP}{\mathbb{P}}
\newcommand{\E}{\mathbb{E}}
\newcommand{\eps}{\varepsilon}
\newcommand{\dd}{\,\mathrm{d}}
\begin{document}
\section*{UnsolvedMath ID: 30003698}
\textbf{Source problem code:} catalogue entry
\subsection{Definitions and mathematical statement}
For each integer \(n\geq 1\), let \(K_n\) be the complete graph on \(\{1,\ldots,n\}\), and let \(B_n=M(K_n)\) be its graphic matroid: its ground set is \(E(K_n)\), and its independent sets are the edge sets containing no cycle. Thus \(B_n\) has rank \(n-1\). We use this explicit indexing convention for the braid matroids; the supplied record does not specify an indexing convention.

To define the invariant, let \(M\) be a finite loopless matroid on a ground set \(E\), with rank function \(r_M\) and total rank \(r=r_M(E)\). A flat is a subset \(F\subseteq E\) such that \(r_M(F\cup\{e\})>r_M(F)\) for every \(e\in E\setminus F\). Write \(L(M)\) for the poset of flats ordered by inclusion, and let \(\mu\) be its M\"obius function, characterized by \(\mu(X,X)=1\) and \(\sum_{X\subseteq Z\subseteq Y}\mu(X,Z)=0\) when \(X\subsetneq Y\), with the sum over flats. Define
\[
\chi_M(t)=\sum_{F\in L(M)}\mu(\varnothing,F)t^{r-r_M(F)}.
\]
For a flat \(F\), the restriction \(M|F\) has rank function \(r_{M|F}(A)=r_M(A)\) for \(A\subseteq F\), and the contraction \(M/F\) has rank function
\[
r_{M/F}(A)=r_M(A\cup F)-r_M(F),\qquad A\subseteq E\setminus F.
\]
The ordinary matroid Kazhdan--Lusztig polynomial \(P_M(t)\in\Z[t]\) is characterized by \(P_M(t)=1\) when \(r=0\), and, when \(r>0\), by
\[
\deg P_M(t)<\frac r2,
\qquad
 t^rP_M(t^{-1})=\sum_{F\in L(M)}\chi_{M|F}(t)P_{M/F}(t).
\]

\textbf{Problem.} Find a simple closed formula, valid for every integer \(n\geq1\), for the entire polynomial \(P_{B_n}(t)\), equivalently for all of its coefficients as functions of \(n\) and the coefficient index.

The phrase ``simple closed formula'' is retained from the source and has no supplied formal criterion. Consequently, the desired form of an answer, and which existing generating or combinatorial descriptions would qualify, remain unresolved. The use of ordinary matroid Kazhdan--Lusztig polynomials is the standard interpretation of the supplied wording; the record does not request an equivariant refinement.
\subsection{Short English statement}
Find a simple closed expression for the full Kazhdan--Lusztig polynomial of the graphic matroid of every complete graph. The source does not define what counts as a simple closed formula. Its supplied, dated assessment labels the problem partially solved and describes progress on top coefficients and computation; that assessment requires independent verification.
\subsection{Sources}
\begin{itemize}
\item Supplied local UnsolvedMath registry record, permanent ID \(30003698\), in the Oberwolfach Reports Open Problems collection; exact source problem code \texttt{catalogue entry}. No individual proposer is identified. The supplied proposed year is \(2018\).
\item Supplied report citation: \emph{Topology of Arrangements and Representation Stability} (2019), Oberwolfach Reports, DOI \texttt{10.4171/owr/2018/2}, pp.~43--123. \url{https://doi.org/10.4171/owr/2018/2}. This reference is supplied as the source attribution; its contents have not been independently checked here.
\item Supplied literature reference: Ferroni and Larson, ``braid matroid polynomials.'' \url{https://arxiv.org/abs/2303.02253}. The record supplies no fuller bibliographic citation.
\item Supplied literature reference: ``Oberwolfach report.'' \url{https://doi.org/10.4171/owr/2018/2}.
\item Literature assessment checked at \texttt{2026-08-21}, according to the supplied record. This is supplied dated curation, not an independently verified assessment. It labels the problem \texttt{partially\_solved} and states that explicit generating and combinatorial results determine the top coefficient and substantially advance computation, while a simple closed formula for the full braid-matroid polynomials is still not established. The background expressly requires independent verification of time-sensitive status and cited claims.
\end{itemize}
\end{document}
